TOYCO produce y vende tres tipos de juguetes: trenes, camiones y coches. 
Cada juguete requiere un proceso de ensamblaje y uno final de empaquetado. 

Los límites semanales  en los tiempos disponibles para las dos operaciones son \textbf{4300 y 4600}, respectivamente, y los márgenes por tren, camión y coche son de \textbf{2, 3 y 5 }euros, respectivamente. Los tiempos de montaje

Los tiempos que requiere cada tren un cada proceso son \textbf{2, 3} minutos, respectivamente. Los de cada camión son \textbf{2 y 4} y los de cada coche \textbf{4 y 4} minutos.

Adémas, existe un compromiso comercial con un distribuidor de juguetes de suministrarle al menos \textbf{1000} unidades entre camiones y coches.

Si $x_1$, $x_2$ y $x_3$ representan el número semanal de unidades producidas de trenes, camiones y coches, respectivamente, el modelo de programación lineal que permite establecer la producción con mayor margen es el siguiente.


\begin{equation}
\begin{split}
\mbox{max. } z = 2x_1 + 3x_2 + 5x_3\\
s.a.:\\
2x_1 + 2x_2 + 4x_3 \leq 4300\\
3x_1 + 4x_2 + 4x_3 \leq 4600\\
        x_2 +  x_3 \geq 1000\\
x_1,\,\,x_2,\,\,x_3\geq 0\\
\end{split}
\end{equation}

La tabla del método del simplex correspondiente a la solución óptima es:

\begin{center}
\begin{tabular}{c|c|cccccc|}
 & $z$ & $x_1$ & $x_2$ & $x_3$ & $h_1$ & $h_2$ & $h_3$\\ 
      & -5450 & $-3/4$ & 0 & 0 & -1 & $-1/4$ & 0 \\ 
\hline
$x_3$ & 1000  & $1/4$ & 0 & 1 & $1/2$ & $-1/4$ & 0\\ 
$h_3$ & 150  & $3/4$ & 0 & 0 & 0 & $1/4$ & 1\\ 
$x_2$ & 150  & $1/2$ & 1 & 0 & $-1/2$ & $1/2$ & 0\\ 
\hline
\end{tabular}
\end{center}





Se pide:

\begin{enumerate}
    \item Interpretar la información que ofrece a Toyco la tabla de la solución óptima.

    De acuerdo con los valores de las variables en la solución óptima:
    \begin{itemize}
        \item No se produce ningún tren.
        \item Se producen 150 camiones.
        \item Se producen 1000 coches.
        \item Se hace uso de todos los minutos de montaje.
        \item Se hace uso de todos los minutos de empaquetado.
        \item Se sirven 150 unidades por encima del mínimo del compromiso comercial adquirido.
        \item El margen total semanal es de 5450 euros.
    \end{itemize}

    \item Dentro de qué rango de valores del tiempo disponible de empaquetado la gama de productos del plan óptima de producción no cambiaría.

    \begin{equation}
    \begin{split}
    u^B=B^{-1}b = \begin{pmatrix}
    \frac{1}{2} & \frac{-1}{4} & 0\\
    0 & \frac{1}{4} & -1\\
    \frac{-1}{2} & \frac{1}{2} & 0\\
    \end{pmatrix}
    \begin{pmatrix}
    4300\\
    b_2\\
    1000\\
    \end{pmatrix}
     = \begin{pmatrix}
    2150 - \frac{b_2}{4}\\
    \frac{b_2}{4} - 1000\\
    -2150 + \frac{b_2}{2}\\
    \end{pmatrix}
    \geq 0 \Rightarrow \\
    4300 \leq b_2 \leq 8600
    \end{split}
    \end{equation}

    Siempre y cuando se disponga de una cantidad de minutos de montaje igual o superior a 4300 y menor o igual que 8500 la gama de productos sería la misma.

    \item ¿Cuál sería el impacto de relajar el compromiso comercial adquirido?

    El precio sombra de la tercera restricción es 0 ($\pi_3=V^B_{h_3}$, por ser restricción de tipo mayor o igual), con lo que aumentar o reducir entorno al compromiso actual de entregar 100 unidades no cambiaría el plan de producción. De hecho, Toyco está entregando 150 unidades por encima de la cantidad acordada.

    \item Por política de empresa, hay que fabricar al menos tantos trenes como camiones, ¿como afectaría esto al plan óptimo de producción?

    Este nuevo requisito daría lugar a una nueva restricción: $x_1 \geq x_2$, equivalente a $x_1 - x_2 \geq 0$ o, como restricción de igualdad, $x_1 - x_2 - h_4 = 0$.


    \begin{center}
    \begin{tabular}{c|c|ccccccc|}
      & $z$ & $x_1$ & $x_2$ & $x_3$ & $h_1$ & $h_2$ & $h_3$ & $h_3$\\ 
      & -5450 & -3/4 & 0 & 0 & -1 & -1/4 & 0 & 0\\ 
    \hline
    $x_3$ & 1000  & $\frac{1}{4}$ & 0 & 1 & $\frac{1}{2}$ & $\frac{-1}{4}$ & 0 & 0 \\ 
    $h_3$ & 150  & $\frac{3}{4}$ & 0 & 0 & 0 & $\frac{1}{4}$ & 1 & 0\\ 
    $x_2$ & 150  & $\frac{1}{2}$ & 1 & 0 & $\frac{-1}{2}$ & $\frac{1}{2}$ & 0 & 0\\ 
    \hline
          & 0     & 1          & -1  & 0 & 0 & 0 & 0 & -1\\           
    \hline
    $h_4$ & -150   & $-\frac{3}{2}$ & 0  & 0 & $\frac{1}{2}$ & $-\frac{1}{2}$ & 0 & 1\\ 
    \hline
       & -5375 & 0 & 0 & 0 & $\frac{-5}{4}$ & 0 & 0 & $\frac{-1}{2}$ \\ 
    \hline
    $x_3$ & 975  & 0 & 0 & 1 & $\frac{7}{12}$ & $\frac{-1}{3}$ & 0 & $\frac{1}{6}$\\ 
    $h_3$ & 75  & 0 & 0 & 0 & $\frac{1}{4}$ & 0 & 1 & $\frac{1}{2}$\\ 
    $x_2$ & 100  & 0 & 1 & 0 & $\frac{-1}{3}$ & $\frac{1}{3}$ & 0 & $\frac{1}{3}$\\ 
    $x_1$ & 100  & 1 & 0 & 0 & $\frac{-1}{3}$ & $\frac{1}{3}$ & 0 & $\frac{-2}{3}$\\ 
    \hline
      & -5375 & 0 & 0 & 0 & $\frac{-5}{4}$ & 0 & 0 & $\frac{-1}{2}$ \\ 
    \hline
    $x_3$ & 1075  & 1 & 0 & 1 & $\frac{1}{4}$ & 0 & 0 & $\frac{-1}{2}$\\ 
    $h_3$ & 75  & 0 & 0 & 0 & $\frac{1}{4}$ & 0 & 1 & $\frac{1}{2}$\\ 
    $x_2$ & 0  & -1 & 1 & 0 & 0 & 0 & 0 & 1\\ 
    $h_2$ & 300  & 3 & 0 & 0 & -1 & 1 & 0 & -2\\ 
    \hline
    \end{tabular}
    \end{center}

    La solución óptima no cumple con la nueva restricción, lo cual hace que la solución básica correspondiente no sea factible ($h_4 = -150 \leq 0$). Aplicando Lemke se puede llegar a dos posibles soluciones alternativas. En ambas el margen total es de 5375 euros. Además:
    \begin{itemize}
        \item En una de ellas se fabrican 100 trenes, 100 camiones y 975 coches.
        \item En la otra solución se fabrican únicamente 1075 coches (no se fabrican trenes ni coches, con lo que el requisito comercial se cumple).
    \end{itemize}



    
\end{enumerate}
